\scanpage{457}\mspage{369}
% archive sheet number: 456
$K$ de but $e$, on a
\[
p'u'v=p'\mathrm{id}_e=p'
\]
donc le composé en question s'écrit aussi
\diagram{\begin{tikzcd}[ampersand replacement=\&]
k\arrow[r,"u"] \& e\arrow[r,"p'p"] \& e\arrow[r,"v'"] \& k'',
\end{tikzcd}}
où $p'p\in\Sigma_\infty$ puisque $\Sigma_\infty$ est stable par composition, donc il est dans \deleted{\unsure{$K'$}} $F$, qed.

Il est clair que $K'$ contient $K$, et ayant $\unsure{=\mathrm{s}|\Rightarrow}$ objets que $K$ qui admet $e$ comme plus grd élémt, que $e$ est un plus grd \unsure{élémt|élément} de $K'$. Reste à prouver que $K'$ est filtrante, Comme $K'$ a un plus grd élémt, il reste à voir \deleted{\illegible{1}} que toute double-flèche \inserted{$k\rightrightarrows l$} de $K'$ peut être égalisée dans $K'$. Quitte à composer avec $l\longrightarrow e$, \unsure{OPS.|OPL.} $l=e$. Or on \unsure{sait|voit} que toute flèche $k\longrightarrow e$ dans $K'$ est de la forme
\diagram{\begin{tikzcd}[ampersand replacement=\&]
k\arrow[r,"u"] \& e\arrow[r,"p"] \& e\arrow[r,"v"] \& \deleted{\unsure{k}}e,
\end{tikzcd}}
\deleted{et quitte à composer} avec $u,v\in\mathrm{Fl}\,K$ et $p\in\Sigma_\infty\cup\{\mathrm{id}_e\}$. Quitte à composer avec $p_0\in\Sigma_\infty$ \deleted{\unsure{choisi|fixé}}\unsure{fixé|choisi}, cette \deleted{\unsure{double}} flèche devient (puisque $p_0v=p_0\mathrm{id}_e=p_0$)
\diagram{\begin{tikzcd}[ampersand replacement=\&]
k\arrow[r,"u"] \& e\arrow[r,"p_0p"] \& e\unsure{..|:}
\end{tikzcd}}
\deleted{ou encore} i.e. $p_0pu$ \deleted{\unsure{pour notre}} Pour une flèche $v'p'u':k\longrightarrow e$, on \unsure{trouvera|trouve} $p_0p'u'$. Soit alors $q\in\Sigma_0$ tel que $q(p_0p)=q(p_0p')$, \inserted{soit $r$} ($q$\unsure{.} existe puisque $p_0p,p_0p'\in\Sigma_\infty$). Donc
\[
\big(q(p_0pu),q(\unsure{p_0^{.}|p_0'|p_0}p'u')\unsure{.}\big)=(r\unsure{u},ru'),
\]
et on a $ru=ru'$ puisque $r\in\Sigma_\infty$. Au total, $qp_0\deleted{\unsure{q}}$ égalise les deux flèches $vpu$ et $v'p'u'$, qed.
